%==============================================================================
% GianniLambrecht-BPvoorstel.tex
% HOOFDBESTAND voor onderzoeksvoorstel — gescheiden inhoud (voorstel-inhoud.tex)
%==============================================================================

\documentclass{hogent-article}

% Invoegen bibliografiebestand
\addbibresource{voorstel.bib}

% Informatie over de opleiding, het vak en soort opdracht
\studyprogramme{Professionele bachelor toegepaste informatica}
\course{Bachelorproef}
\assignmenttype{Onderzoeksvoorstel}

\academicyear{2025-2026} % pas aan indien nodig

% Werktitel
\title{Een gebruiksvriendelijke planningsapplicatie op maat van mensen met ADHD: ontwerp, vergelijking en optimalisatie van digitale hulpmiddelen voor beter taakbeheer}

% Studentnaam en emailadres
\author{Gianni Lambrecht}
\email{gianni.lambrecht@student.hogent.be}

% Geen co-promotor voor nu
% \supervisor[Co-promotor]{S. Beekman (Synalco, \href{mailto:sigrid.beekman@synalco.be}{sigrid.beekman@synalco.be})}

% Specialisatierichting
\specialisation{Mobile \& Enterprise development}
\keywords{ADHD, planningsapp, UX, toegankelijkheid, gamification}

\begin{document}

\begin{abstract}
Een korte samenvatting van het voorstel: dit onderzoek richt zich op het ontwerp en de evaluatie van een digitale planningsapplicatie die specifiek inspeelt op de noden van personen met ADHD. Uit eerdere analyse blijkt dat bestaande planningsapps vaak te complex of overweldigend zijn en onvoldoende automatische ondersteuning bieden bij het opdelen van taken. Het doel is een wetenschappelijk onderbouwd prototype te ontwikkelen en te vergelijken met gangbare tools op gebruiksvriendelijkheid en effect op uitstelgedrag en ervaren stress. De voorgestelde methodologie omvat literatuurstudie, behoefteanalyse met gebruikers (interviews / korte gebruikerstesten), ontwerp en iteratieve ontwikkeling van een proof-of-concept en een vergelijkende evaluatie. Verwachte resultaten zijn een set ontwerpprincipes, een werkend prototype en concrete aanbevelingen voor de ontwikkeling van ADHD-vriendelijke planningshulpmiddelen.
\end{abstract}

\tableofcontents

% De hoofdtekst (gescheiden bestand)
%==============================================================================
% voorstel-inhoud.tex
% Inhoud voor het onderzoeksvoorstel van Gianni Lambrecht
%==============================================================================

%---------- Inleiding ---------------------------------------------------------
\section{Introductie}
\label{sec:inleiding}

Waarover zal je bachelorproef gaan? Introduceer het thema en zorg dat volgende zaken zeker duidelijk aanwezig zijn:
\begin{itemize}
  \item kaderen thema
  \item de doelgroep
  \item de probleemstelling en (centrale) onderzoeksvraag
  \item de onderzoeksdoelstelling
\end{itemize}

Personen met ADHD (Attention Deficit Hyperactivity Disorder) ervaren vaak moeilijkheden met executieve functies zoals planning, organisatie, tijdsbeheer en taakvoltooiing. Veel bestaande planningsapps (bijv. Todoist, Notion, TickTick, Tiimo, Trello) richten zich op brede gebruikersgroepen en missen vaak specifieke aanpassingen die nodig zijn voor ADHD-gebruikers. Deze apps kunnen te complex zijn, onvoldoende structuur bieden, of notificaties gebruiken die extra stress veroorzaken in plaats van te helpen. Daardoor is er behoefte aan een toepassing die rekening houdt met beperkte werkgeheugencapaciteit, neiging tot uitstel en gevoeligheid voor prikkels.

\medskip
\textbf{Doelgroep:} jongvolwassenen en volwassenen met ADHD die digitale hulpmiddelen gebruiken voor taakbeheer, studie of werk.

\medskip
\textbf{Casus en afbakening:} dit onderzoek richt zich op het ontwerpen, implementeren en evalueren van een prototype van een planningsapp die inspeelt op de noden van ADHD-gebruikers. De focus ligt op UX-ontwerp, automatische taakopsplitsing, rustige notificaties en motiverende, niet-overprikkelende feedbackmechanismen.

\medskip
\textbf{Probleemstelling:} hoewel er veel planningsapplicaties bestaan, blijken deze tools voor veel ADHD-gebruikers te complex en niet afgestemd op hun specifieke cognitieve en emotionele noden. Er is nood aan een wetenschappelijk onderbouwde analyse en een prototype van een tool die automatische taakverdeling, aangepaste UX en stressreductie combineert.

\medskip
\textbf{Centrale onderzoeksvraag:}
\begin{quote}
Hoe kan een applicatie op maat van mensen met ADHD bijdragen aan beter taakbeheer en tijdsplanning, met minder stress en uitstelgedrag?
\end{quote}

\medskip
\textbf{Onderzoeksdoelstelling:} het ontwikkelen en evalueren van een proof-of-concept dat ontwerpprincipes en technische realisaties (front-end en eenvoudige backend) demonstreert, aangevuld met aanbevelingen voor opschaling en verdere validatie.

%---------- State-of-the-art -------------------------------------------------
\section{State-of-the-art}
\label{sec:literatuurstudie}

Hier beschrijf je de state-of-the-art rondom je gekozen onderzoeksdomein. Je steunt daarbij sterk op professionele vakliteratuur en recente studies.

ADHD beïnvloedt vooral executieve functies zoals werkgeheugen en planningcapaciteit; dit vertaalt zich in moeilijkheden bij het opdelen en volhouden van taken en in een afwijkende tijdsbeleving. Studies tonen dat visuele structuur, eenvoudige workflows en beperkte prikkelintensiteit kunnen helpen om taakvoltooiing te verbeteren. Daarnaast wijzen onderzoeken op de rol van beloningsmechanismen en aangepaste notificaties om motivatie te stimuleren zonder extra stress te veroorzaken.

Bestaande apps (bijv. Tiimo) bieden interessante voorbeelden van visuele routines, maar tonen beperkingen in flexibiliteit en automatisering voor complexe taken. Er is een lacune voor apps die:
\begin{itemize}
  \item automatische opsplitsing van taken in haalbare stappen bieden;
  \item aanpasbare visuele beloningen tonen zonder overprikkeling;
  \item rustige en contextbewuste notificaties hanteren;
  \item eenvoudige, voorspelbare UI-flows hebben.
\end{itemize}

In de literatuur zal ik werken met peer-reviewed bronnen over ADHD en executieve functies, HCI-onderzoek over cognitieve toegankelijkheid, en empirische evaluaties van bestaande apps en prototypes. Verwijzingen naar gebruikte bronnen vind je in de Referenties.

%---------- Methodologie ------------------------------------------------------
\section{Methodologie}
\label{sec:methodologie}

Hier beschrijf je hoe je van plan bent het onderzoek te voeren. Welke onderzoekstechniek ga je toepassen om elk van je onderzoeksvragen te beantwoorden?

De voorgestelde methodologie bestaat uit de volgende fasen:

\begin{enumerate}
  \item \textbf{Literatuurstudie}: systematische verkenning van wetenschappelijke publicaties en relevante gray literature omtrent ADHD, UX-ontwerp voor cognitieve toegankelijkheid en bestaande planningsapps.
  \item \textbf{Behoefteanalyse (kwalitatief)}: semigestructureerde interviews en korte gebruikerstesten (5–10 deelnemers met ADHD) om concrete pijnpunten en wensen te identificeren. Deze fase genereert requirements en prioriteiten.
  \item \textbf{Ontwerp}: opstellen van ontwerpprincipes, wireframes en interactieve mock-ups (Figma). Iteratieve feedbacksessies met 2–3 gebruikers.
  \item \textbf{Implementatie (PoC)}: bouw van een eenvoudig maar functioneel prototype (React of React Native; backend met Firebase of Supabase) met nadruk op automatische taakopsplitsing, visuele feedback en contextbewuste notificaties.
  \item \textbf{Evaluatie}: vergelijkende gebruikerstesten (within-subjects of matched) tussen prototype en bestaande app(s). Meetinstrumenten: SUS (usability), korte stress-vragenlijst (zelfgerapporteerd), taakvoltooiingspercentage en kwalitatieve feedback.
  \item \textbf{Analyse en synthese}: kwantitatieve analyse van meetgegevens en thematische analyse van kwalitatieve data. Formuleren van aanbevelingen en discussie over beperkingen.
\end{enumerate}

\medskip
\textbf{Tools en technologieën:}
\begin{itemize}
  \item Ontwerp: Figma voor wireframes en prototyping.
  \item Implementatie: React (web) of React Native (mobiel); backend: Firebase/Supabase.
  \item Analyse: Python (pandas, scipy) of R voor eenvoudige statistische tests; NVivo/Atlas.ti of handmatige thematische codering voor kwalitatieve data.
\end{itemize}

\medskip
\textbf{Tijdsplanning en deliverables (globaal):}
\begin{itemize}
  \item Maand 1–2: literatuurstudie en opzet methodologie (deliverable: literatuurmatrix).
  \item Maand 3: behoefteanalyse en requirements (deliverable: requirementsdoc).
  \item Maand 4–5: ontwerp en prototypeontwikkeling (deliverable: Figma-mockups + PoC).
  \item Maand 6: gebruikerstesten en gegevensverzameling (deliverable: testverslag).
  \item Maand 7: analyse, schrijven en aanbevelingen (deliverable: scriptie + codebase).
\end{itemize}

%---------- Verwacht resultaat, conclusie ------------------------------------
\section{Verwacht resultaat, conclusie}
\label{sec:verwachte_resultaten}

Hier beschrijf je welke resultaten je verwacht.

Verwacht wordt:
\begin{itemize}
  \item een overzicht van pijnpunten en requirements voor ADHD-vriendelijke planningstools;
  \item een set ontwerpprincipes en richtlijnen voor cognitief toegankelijke planningsapps;
  \item een werkend prototype dat aantoont hoe automatische taakopsplitsing en rustige notificaties in de praktijk kunnen werken;
  \item empirische resultaten die aantonen of gebruikers het prototype als gebruiksvriendelijker en minder stressvol ervaren dan bestaande oplossingen.
\end{itemize}

De meerwaarde: praktische richtlijnen en een concrete technische demonstrator voor ontwikkelaars, coaches en hulpverleners die digitale hulpmiddelen willen inzetten voor mensen met ADHD.



\nocite{*}
\printbibliography[heading=bibintoc,title={Referenties}]

\end{document}
